\begin{table}[t]
\centering
\caption{Fresh-state query workloads at a common cost tolerance}
\label{tab:r5_workload}
\small
\begin{tabular}{rrrrrr}
\toprule
Queries & Max. loss bound & Training & Query & Reference & Cost ratio \\
\midrule
32 & $6.024\times10^{-4}$ & 0.812 & 0.040 & 0.152 & 0.18 \\
128 & $6.574\times10^{-4}$ & 0.812 & 0.119 & 0.420 & 0.45 \\
512 & $6.574\times10^{-4}$ & 0.812 & 0.600 & 1.477 & 1.05 \\
2048 & $6.754\times10^{-4}$ & 0.812 & 2.600 & 5.708 & 1.67 \\
\bottomrule
\end{tabular}
\par\medskip
\begin{minipage}{.97\linewidth}
\footnotesize Dimension four, seed 101; held-out state seed 541021, uniform on $[0,1.2]^4$. Larger workloads extend the same fresh-state sequence. Times are seconds. Query time includes zero-start convex improvement and all 64 terminal paths per initial state. Reference time is for independent nonanticipative optimization to the common $10^{-3}$ cost tolerance. The cost ratio is reference time divided by critic training plus query time. Each reported maximum policy-loss bound is below $10^{-3}$. Timings are medians of three evaluations; the separate tighter validation reference is not charged to either competing solver.
\end{minipage}
\end{table}
